% §7.3 Confidence Intervals Properly
%
% Cross-reference resolution status:
%
% DEFINED here:
%   \label{sec:confidence-intervals}              -- this section
%   \label{def:wald-ci}                           -- Wald CI definition
%   \label{def:wilson-ci}                         -- Wilson CI definition
%   \label{def:clopper-pearson-ci}                -- Clopper-Pearson CI definition
%   \label{def:jeffreys-ci}                       -- Jeffreys CI definition
%
% ALREADY DEFINED (chapter 1):
%   \Cref{sec:sampling-rates}                     -- chapter 1 §1.4
%
% ALREADY DEFINED (chapter 7, this plan):
%   \Cref{sec:setup-revisited}                    -- §7.1
%   \Cref{sec:binomial-estimator}                 -- §7.2
%   \Cref{prop:binomial-variance}                 -- §7.2 variance proposition
%
% FORWARD REFS within chapter 7 (resolve in later Plan 7 tasks):
%   \Cref{sec:derived-measures}                   -- §7.4 (Task 5)
%
% §7.3 retires items (ii) and (iii) of the §1.4 deferral named in §7.1:
% the proper-CI question (Wald, Wilson, Clopper-Pearson, Jeffreys) and
% the sawtooth under-coverage observed in the chapter 1 notebook. The
% four CI definitions cumulate; the closing recommendation is Wilson by
% default, Jeffreys as the Bayesian alternative, Clopper-Pearson when
% conservative coverage is required, Wald only in large-$n_0$ regimes
% away from the boundaries.

\section{Confidence Intervals Properly}
\label{sec:confidence-intervals}

\Cref{sec:binomial-estimator} fixed the point estimate
$\hat{\fprate} = \mathrm{FP}/n_0$ as the unbiased, efficient
maximum-likelihood estimator of $\fprate$. A point estimate is not a
confidence statement. The Wald interval given in
\Cref{sec:sampling-rates} attaches Normal-approximation error bars to
the point estimate and stops there. The chapter~1 notebook showed the
attached error bars failing in exactly the regime a practitioner most
needs them: small $n_0$, $\fprate$ near $0$ or $1$. This section
diagnoses the failure rigorously, then derives three alternatives
(Wilson, Clopper-Pearson, Jeffreys), compares their coverage on the
small-$n_0$ regime, and recommends a default.

\subsection*{The Wald interval}

The Wald interval is the central-limit-theorem interval read off the
sampling distribution of the point estimate, with the unknown
$\fprate$ in the variance replaced by its empirical value.

\begin{definition}[Wald confidence interval]
\label{def:wald-ci}
At confidence level $1 - \alpha$, the Wald interval for $\fprate$ is
\[
  \hat{\fprate} \;\pm\; z_{\alpha/2}
    \sqrt{\frac{\hat{\fprate}(1 - \hat{\fprate})}{n_0}},
\]
where $z_{\alpha/2}$ is the upper-$\alpha/2$ quantile of the standard
Normal distribution and $\hat{\fprate} = \mathrm{FP}/n_0$ is the
point estimate from \Cref{sec:binomial-estimator}.
\end{definition}

The derivation is the standard one. The central limit theorem applied
to the i.i.d.\ sum $\mathrm{FP} = \sum X_i$ of
\Cref{sec:binomial-estimator} gives an asymptotic Normal approximation
$\hat{\fprate} \sim \mathrm{Normal}(\fprate, \fprate(1 - \fprate)/n_0)$
as $n_0 \to \infty$. Substituting $\hat{\fprate}$ for $\fprate$ inside
the variance (the plug-in step) and reading off the $1 - \alpha$
two-sided quantile interval of the resulting Normal yields the
formula. The interval is correct in the limit $n_0 \to \infty$ for
every fixed $\fprate \in (0, 1)$. The trouble is that practitioners
work at finite $n_0$.

\subsection*{The under-coverage problem}

The chapter~1 notebook ran $M$ independent replications of the
classifier at sample size $n_0$, constructed the Wald interval at
each, and tallied the fraction containing the true $\fprate$. At
$n_0 = 100$ and $\fprate = 0.05$ the empirical coverage settled near
$0.92$ rather than the nominal $0.95$. The gap widened as $\fprate$
moved closer to $0$. Two distinct failure modes are at work.

\smallskip
\noindent\emph{(a) The boundary problem.}
The Wald formula is a symmetric interval centered at $\hat{\fprate}$,
which lives in $[0, 1]$. The radius
$z_{\alpha/2} \sqrt{\hat{\fprate}(1 - \hat{\fprate})/n_0}$ is positive
and depends on $\hat{\fprate}$, not on $\fprate$. When
$\hat{\fprate}$ is small and the radius is comparable to it, the
lower endpoint $\hat{\fprate} - z_{\alpha/2}
\sqrt{\hat{\fprate}(1 - \hat{\fprate})/n_0}$ can fall below $0$;
symmetrically the upper endpoint can rise above $1$ when
$\hat{\fprate}$ is near $1$. Both cases assign positive probability
mass to impossible values of $\fprate$. The interval has to be
clipped to $[0, 1]$ by hand, which is not a property the construction
provides.

\smallskip
\noindent\emph{(b) The skew problem.}
The binomial distribution $\mathrm{Binomial}(n_0, \fprate)$ is
symmetric only at $\fprate = 1/2$. Near $\fprate = 0$ it is
right-skewed: the mass piles up at small counts with a long right
tail. Near $\fprate = 1$ it is left-skewed. The Normal approximation
is symmetric and cannot capture the skew. At small $n_0$ and
$\fprate$ near a boundary, the central limit theorem has not yet
overcome the binomial skew, and the symmetric Wald interval
miscalibrates: it places too much of its tail probability on the
side away from the boundary and too little on the side toward it.
The empirical coverage drops below nominal.

\smallskip
There is a third pattern, visible in the chapter~1 notebook sweep of
$n_0$ at fixed $\fprate = 0.05$. The empirical coverage as a
function of $n_0$ is not monotone: it rises and dips in a jagged
pattern as $n_0$ grows. This is the \emph{sawtooth}. The cause is
discreteness. The Wald interval is built on a continuous Normal
approximation, but the underlying $\mathrm{FP}$ takes only the
integer values $0, 1, \dots, n_0$. As $n_0$ increments by one, the
boundary at which the interval just barely covers or just barely
misses the true $\fprate$ shifts in a discrete jump, and the
empirical coverage hops between values close to but rarely exactly
at nominal. The pattern is a real feature of every CI built on a
discrete distribution; the Wald interval is unusual in how
pronounced the sawtooth is, because of failure modes (a) and (b)
above.

The three alternatives in this section address the failures in
different ways: Wilson fixes the boundary problem by inverting a
different statistic, Clopper-Pearson sidesteps the Normal
approximation entirely by inverting the exact binomial CDF, and
Jeffreys does the Bayesian thing and reports a posterior credible
interval.

\subsection*{The Wilson interval}

The Wald construction plugs $\hat{\fprate}$ into the variance and
inverts the resulting Normal-form statistic. Wilson inverts a
different statistic, the score statistic, which uses the true
$\fprate$ (not its estimate) in the variance.

The score statistic for the binomial mean is
\[
  Z(\fprate)
  \;=\; \frac{\hat{\fprate} - \fprate}{\sqrt{\fprate(1 - \fprate)/n_0}}.
\]
Under the central limit theorem, $Z(\fprate) \to \mathrm{Normal}(0, 1)$
as $n_0 \to \infty$ when $\fprate$ is the true parameter, so the
confidence set at level $1 - \alpha$ is the set of $\fprate$
satisfying $|Z(\fprate)| \leq z_{\alpha/2}$. Squaring both sides and
clearing the denominator gives a quadratic inequality in $\fprate$:
\[
  (\hat{\fprate} - \fprate)^2 \;\leq\; z_{\alpha/2}^2 \,
    \frac{\fprate(1 - \fprate)}{n_0}.
\]
Writing $z = z_{\alpha/2}$ for brevity and expanding,
\[
  \hat{\fprate}^2 - 2 \hat{\fprate} \fprate + \fprate^2
  \;\leq\;
  \frac{z^2}{n_0} \fprate - \frac{z^2}{n_0} \fprate^2,
\]
which rearranges to
\[
  \fprate^2 \left(1 + \frac{z^2}{n_0}\right)
    - 2 \fprate \left(\hat{\fprate} + \frac{z^2}{2 n_0}\right)
    + \hat{\fprate}^2 \;\leq\; 0.
\]
This is a downward-closed quadratic in $\fprate$, so its solution set
is the interval between the two roots. Applying the quadratic formula
and simplifying gives the form below.

\begin{definition}[Wilson confidence interval]
\label{def:wilson-ci}
At confidence level $1 - \alpha$, with $z = z_{\alpha/2}$, the Wilson
interval for $\fprate$ is
\[
  \frac{\hat{\fprate} + z^2/(2 n_0)}{1 + z^2/n_0}
  \;\pm\;
  \frac{z}{1 + z^2/n_0}
    \sqrt{\frac{\hat{\fprate}(1 - \hat{\fprate})}{n_0}
      + \frac{z^2}{4 n_0^2}}.
\]
\end{definition}

The center of the Wilson interval is not $\hat{\fprate}$ but a shrunken
version of it pulled toward $1/2$ by the factor $z^2/(2 n_0)$ in the
numerator and the factor $1 + z^2/n_0$ in the denominator. The radius
likewise has a correction $z^2/(4 n_0^2)$ inside the square root that
keeps the interval well-defined even at $\hat{\fprate} \in \{0, 1\}$.
Both adjustments matter. At $\hat{\fprate} = 0$ the Wald interval
collapses to the degenerate point $[0, 0]$; the Wilson interval
gives $[0, z^2/(n_0 + z^2)]$, a non-empty interval whose upper
endpoint admits the possibility that the true rate is small but
positive. The boundary problem of (a) is gone: both endpoints lie in
$[0, 1]$ for every $\hat{\fprate}$ and every $n_0$, because the
defining inequality has the true $\fprate$ in the variance and is a
proper quadratic in $\fprate$. The skew problem of (b) is partially
addressed: the Wilson center is shifted toward $1/2$, which is the
direction the symmetric Wald interval got wrong near the boundaries.

\subsection*{The Clopper-Pearson interval}

Wilson still relies on a Normal approximation, even if a sharper one.
Clopper-Pearson abandons the approximation and works with the exact
binomial CDF. The construction is the test-inversion construction in
its cleanest form: for each candidate $\fprate$, test whether the
observed $\mathrm{FP}$ is consistent with $\mathrm{Binomial}(n_0,
\fprate)$ at level $\alpha$, and include $\fprate$ in the interval iff
the test fails to reject.

The lower endpoint $L$ is the smallest $\fprate$ such that
$\Pr[\mathrm{Binomial}(n_0, \fprate) \geq \mathrm{FP}] > \alpha/2$;
equivalently, $L$ is the largest $\fprate$ such that
$\Pr[\mathrm{Binomial}(n_0, \fprate) \geq \mathrm{FP}] \leq \alpha/2$.
The upper endpoint $U$ is symmetric: the smallest $\fprate$ such that
$\Pr[\mathrm{Binomial}(n_0, \fprate) \leq \mathrm{FP}] \leq \alpha/2$.

These endpoints have closed forms via the Beta distribution, because
the binomial CDF and the Beta CDF are related by the identity
\[
  \Pr[\mathrm{Binomial}(n_0, \fprate) \geq k]
  \;=\; \Pr[\mathrm{Beta}(k, n_0 - k + 1) \leq \fprate].
\]
The endpoints are therefore Beta quantiles.

\begin{definition}[Clopper-Pearson confidence interval]
\label{def:clopper-pearson-ci}
At confidence level $1 - \alpha$, the Clopper-Pearson interval for
$\fprate$ has endpoints
\[
  L \;=\; F^{-1}_{\mathrm{Beta}(\mathrm{FP},\, n_0 - \mathrm{FP} + 1)}(\alpha/2),
  \qquad
  U \;=\; F^{-1}_{\mathrm{Beta}(\mathrm{FP} + 1,\, n_0 - \mathrm{FP})}(1 - \alpha/2),
\]
where $F^{-1}_{\mathrm{Beta}(a, b)}$ denotes the inverse CDF of the
$\mathrm{Beta}(a, b)$ distribution. At the boundaries: if
$\mathrm{FP} = 0$ set $L = 0$; if $\mathrm{FP} = n_0$ set $U = 1$.
\end{definition}

The interval is called \emph{exact} because the construction uses the
exact binomial CDF rather than a Normal approximation. The label is
slightly misleading: the coverage probability is guaranteed to be at
least $1 - \alpha$ at every $\fprate$, but is strictly greater than
$1 - \alpha$ at most $\fprate$ because the binomial is discrete. The
interval over-covers, sometimes substantially, especially at small
$n_0$. Conservative coverage is the price of the guarantee.

\subsection*{The Jeffreys interval}

The Bayesian alternative reports a credible interval from a posterior
distribution over $\fprate$. The choice of prior matters; the
Jeffreys prior is the one whose density is proportional to the square
root of the Fisher information,
\[
  \pi_J(\fprate)
  \;\propto\; \sqrt{I(\fprate)}
  \;=\; \sqrt{\frac{n_0}{\fprate(1 - \fprate)}}
  \;\propto\; \fprate^{-1/2} (1 - \fprate)^{-1/2}.
\]
This is the kernel of the $\mathrm{Beta}(1/2, 1/2)$ distribution, which
serves as a non-informative reference prior with a frequentist
calibration property: the resulting credible interval has approximate
$1 - \alpha$ coverage to higher order than the Wald interval. Under
this prior the posterior given the observed count $\mathrm{FP}$ is
$\mathrm{Beta}(\mathrm{FP} + 1/2, n_0 - \mathrm{FP} + 1/2)$ by
conjugacy.

\begin{definition}[Jeffreys credible interval]
\label{def:jeffreys-ci}
At credibility level $1 - \alpha$, the Jeffreys interval for
$\fprate$ is the central $1 - \alpha$ quantile interval of the
posterior $\mathrm{Beta}(\mathrm{FP} + 1/2, n_0 - \mathrm{FP} + 1/2)$:
\[
  \bigl[\,
    F^{-1}_{\mathrm{Beta}(\mathrm{FP} + 1/2,\, n_0 - \mathrm{FP} + 1/2)}(\alpha/2),
    \;
    F^{-1}_{\mathrm{Beta}(\mathrm{FP} + 1/2,\, n_0 - \mathrm{FP} + 1/2)}(1 - \alpha/2)
  \,\bigr].
\]
\end{definition}

The Jeffreys interval is structurally close to Clopper-Pearson but
uses different Beta parameters (the half-integer shift) and does not
treat the endpoints asymmetrically. At $\mathrm{FP} = 0$ the
posterior is $\mathrm{Beta}(1/2, n_0 + 1/2)$, which has a
well-defined lower quantile strictly above $0$; the interval is
non-degenerate at the boundaries. The frequentist coverage of the
Jeffreys interval is close to nominal at small $n_0$ without the
heavy conservatism of Clopper-Pearson.

\subsection*{Comparison}

The four intervals can be ranked along several axes: empirical
coverage at a fixed $(n_0, \fprate)$, behavior at the boundary
$\hat{\fprate} = 0$, and structural symmetry about the point
estimate. \Cref{tab:ci-comparison} summarizes.

\begin{table}[ht]
\centering
\begin{tabular}{l|cccc}
\hline
                                & Wald             & Wilson           & Clopper-Pearson  & Jeffreys \\
\hline
Coverage at $n_0 = 100$, $\fprate = 0.05$
                                & $\approx 0.92$   & $\approx 0.94$   & $\geq 0.95$      & $\approx 0.945$ \\
Coverage at $n_0 = 1000$, $\fprate = 0.05$
                                & $\approx 0.95$   & $\approx 0.95$   & $\geq 0.95$      & $\approx 0.95$ \\
Interval at $\hat{\fprate} = 0$
                                & $[0, 0]$         & $[0,\, z^2/(n_0 + z^2)]$
                                                                      & $[0,\, U_{\mathrm{CP}}]$
                                                                                         & $[L_J,\, U_J]$, $L_J > 0$ \\
Symmetric about $\hat{\fprate}$ & yes              & no               & no               & no \\
Endpoints in $[0, 1]$ by construction
                                & no               & yes              & yes              & yes \\
Approximation type              & Normal           & Normal (score)   & exact binomial   & Bayesian posterior \\
\hline
\end{tabular}
\caption{Comparison of four binomial confidence intervals at
nominal $1 - \alpha = 0.95$. Coverage numbers are read from the
chapter~1 notebook sweep at $\fprate = 0.05$ for the Wald row;
Wilson, Clopper-Pearson, and Jeffreys numbers are standard results
from Brown, Cai, and DasGupta (2001) at the same parameters. The
small-$n_0$ row shows the Wald under-coverage diagnosed above and
the substantial recovery from each of the three alternatives. The
large-$n_0$ row confirms that all four intervals agree
asymptotically. The boundary row shows that Wald collapses to a
degenerate point at $\hat{\fprate} = 0$ while the three
alternatives give non-trivial intervals that admit small positive
values of $\fprate$ as compatible with the data.}
\label{tab:ci-comparison}
\end{table}

The coverage entries for the small-$n_0$ row are not constants: each
oscillates along a sawtooth as $n_0$ varies, by the discreteness
argument above. The numbers report the average coverage over a small
window of $n_0$ values near $100$, smoothing the sawtooth. The Wilson
and Jeffreys sawteeth are mild and stay close to nominal. The
Clopper-Pearson sawtooth dips down only to the guarantee floor
$1 - \alpha = 0.95$ and is otherwise above; the over-coverage is
visible in the strict inequality. The Wald sawtooth dips below
nominal repeatedly and is the only one of the four whose minimum is
far enough below nominal to mislead at the precision a careful
practitioner would report.

\subsection*{Recommendation}

The Wilson interval is the default for binomial confidence intervals
in this book and is what \Cref{sec:derived-measures} feeds into the
delta-method derivations. It corrects the Wald boundary failure
without the conservatism of Clopper-Pearson and is easy to implement
from a Normal-quantile table. Jeffreys is the Bayesian alternative
when a credible-interval interpretation is wanted; its frequentist
coverage is also excellent. Clopper-Pearson is the conservative-safe
choice when an absolute lower bound on coverage is required (in a
regulatory or worst-case setting) and a slightly wider interval is
acceptable. Wald is acceptable only when $n_0$ is large and
$\hat{\fprate}$ is far from $0$ and $1$, the regime in which all
four intervals agree; using Wald outside that regime invites the
under-coverage of \Cref{tab:ci-comparison}.

\subsection*{What has been retired}

Items~(ii) and~(iii) of the §1.4 deferral named in
\Cref{sec:setup-revisited} are now retired. The Wilson,
Clopper-Pearson, and Jeffreys intervals fix the boundary failure of
the Wald interval and substantially recover the lost coverage in the
small-$n_0$ regime; the sawtooth pattern observed in the chapter~1
notebook is now attributed to its proper cause (binomial
discreteness viewed through a continuous-distribution lens) and shown
to be mild for the three alternatives. \Cref{sec:derived-measures}
takes up item~(iv): the derived measures (PPV, NPV, F1) and the
delta-method translation of binomial CIs on $\fprate, \fnrate$ into
CIs on the derived quantities.
